% Progreso y brecha (~1 min), en el orden del informe (report_05_final). Fuentes: delivery.yaml, EXP-021.
\begin{frame}{Progreso y brecha}
  \begin{columns}[c,onlytextwidth]
    \begin{column}{0.49\textwidth}
      \begin{figure}
        \centering
        \safegraphics[width=\textwidth,height=0.5\textheight,keepaspectratio]{figures/progress.png}{Progreso por paso}
        \caption{\footnotesize \figcaptionshort{progress}}
      \end{figure}
    \end{column}
    \begin{column}{0.49\textwidth}
      \begin{figure}
        \centering
        \safegraphics[width=\textwidth,height=0.5\textheight,keepaspectratio]{figures/overfitting_curves.png}{Curvas de train y validación}
        \caption{\footnotesize \figcaptionshort{overfitting_curves}}
      \end{figure}
    \end{column}
  \end{columns}
  \keymessage{Mean recall \res{baseline-mean-recall} $\rightarrow$ \res{best-mean-recall}; brecha \res{phase0-gap-acc} $\rightarrow$ \res{final-gap-acc}~pp: se reduce, pero no desaparece.}
  \speakernote{1}{BatchNorm, rotaciones y bajar la lr más tarde (+\res{gain-rlr8}~pp; se explica en la diapositiva siguiente). El modelo regularizado aprende más despacio y supera al de la entrega 1 cuando este ya ha parado. Con las mismas paciencias, la brecha llegaba a \res{phase2-gap-acc}~pp; entrenar más recupera algo de memorización a cambio de mejor validación.}
\end{frame}
